%% ma: L10-B15-0083
%% lop: 10
%% chuong: 6
%% bai: 15
%% dang: 10.15.4
%% loai: DS
%% muc: [TH, NB, VD, TH]
%% dap_an: "ĐSĐĐ"
%% nguon: "Ảnh giáo viên gửi 25/11/2026: Bài tập trắc nghiệm – Toán 10 Bài 15 Hàm số (3 phần), Phần II Câu 5"
%% kien_thuc: "Lớp 10 Bài 15 (tập xác định, đồng biến nghịch biến, điểm thuộc đồ thị); hàm số bậc nhất, hằng đẳng thức (lớp 8–9)"
%% trang_thai: chinh-thuc
%% ghi_chu: "Câu tổng hợp nhiều nội dung Bài 15, xếp theo ý b (đồng biến, nghịch biến)."
\begin{ex}
Cho hai hàm số $f(x)=(m-2)x+1$ và $g(x)=\dfrac{2}{x^2-5x+9}$. Xét tính đúng sai của các khẳng định sau:
\choiceTF
{\True Hàm số $g(x)$ có tập xác định $\mathbb R$.}
{Khi $m>2$ hàm số $f(x)$ nghịch biến trên $\mathbb R$.}
{\True Giá trị lớn nhất của hàm số $g(x)$ bằng $\dfrac{8}{11}$.}
{\True Khi $m=\dfrac{11}{6}$ thì đồ thị hàm số của $f(x)$ và $g(x)$ cùng đi qua điểm $\left(2;\dfrac23\right)$.}
\loigiai{a) $x^2-5x+9=\left(x-\tfrac52\right)^2+\tfrac{11}{4}>0$: Đúng. b) $m-2>0$ nên $f$ đồng biến: Sai. c) $x^2-5x+9\ge\tfrac{11}4$ nên $g(x)\le\tfrac8{11}$, dấu bằng tại $x=\tfrac52$: Đúng. d) $f(2)=2\left(\tfrac{11}6-2\right)+1=\tfrac23$, $g(2)=\tfrac23$: Đúng.}
\end{ex}
